\section*{第 \ref{ch:b3:cytoskeleton-motility} 章——细胞骨架动力学与细胞运动}

\begin{solution}{exo:b3:cytoskeleton-motility:1}
肌动蛋白：\qty{7}{nm}，球状肌动蛋白，有极性（带刺端/尖端），ATP；
皮层、突出、\omterm{def:b3:cytoskeleton-motility:motors}{肌球蛋白}介导的收缩。\omterm{def:b3:cytoskeleton-motility:filaments}{微管}：\qty{25}{nm}，
$\alpha\beta$-微管蛋白二聚体，有极性（正端/负端），GTP；远距离
运输、纺锤体、\omterm{def:b3:cytoskeleton-motility:cilia}{纤毛}。\omterm{def:b3:cytoskeleton-motility:filaments}{中间纤维}：\qty{10}{nm}，
卷曲螺旋蛋白（角蛋白、波形蛋白、核纤层蛋白），无极性，不含核苷酸；
维持细胞和细胞核的机械强度。
\end{solution}

\begin{solution}{exo:b3:cytoskeleton-motility:2}
使某一端既不生长也不缩短的单体浓度，
即 $k_{\mathrm{off}}/k_{\mathrm{on}}$。两端不同，是由于亚基
在掺入后水解自身的 ATP，因此两端呈现不同的
化学物种（到达带刺端的 ATP-肌动蛋白，
从尖端离开的 ADP-肌动蛋白），亲和力不同。若两端\omterm{thm:b3:cytoskeleton-motility:treadmilling}{临界
浓度}相同，则只会有单一平衡：不会出现\omterm{thm:b3:cytoskeleton-motility:treadmilling}{踏车行为}，
没有通量，也无法为定向周转提供能量。
\end{solution}

\begin{solution}{exo:b3:cytoskeleton-motility:3}
外周（\omterm{def:b3:cytoskeleton-motility:filaments}{微管}上）：\omterm{def:b3:cytoskeleton-motility:motors}{驱动蛋白}，朝向正端。中心：
胞质\omterm{def:b3:cytoskeleton-motility:motors}{动力蛋白}，朝向负端。应力纤维：\omterm{def:b3:cytoskeleton-motility:motors}{肌球蛋白} II，
朝向肌动蛋白的带刺端。\omterm{def:b3:cytoskeleton-motility:cilia}{纤毛}：\omterm{def:b3:cytoskeleton-motility:cilia}{轴丝动力蛋白}，朝向
相邻二联管的负端（基部）。
\end{solution}

\begin{solution}{exo:b3:cytoskeleton-motility:4}
突出（\omterm{def:b3:cytoskeleton-motility:migration}{片状伪足}由 Rac 调控，丝状伪足由 Cdc42 调控），黏附
（\omterm{def:b3:cytoskeleton-motility:migration}{整合素}；位于 Rac 和 Rho 的下游），收缩（Rho，通过
\omterm{def:b3:cytoskeleton-motility:motors}{肌球蛋白} II），后部回缩（黏附释放，由
收缩驱动）。
\end{solution}

\begin{solution}{exo:b3:cytoskeleton-motility:5}
$C_{c}^{+} = 1/5 = \qty{0.2}{\micro M}$, $C_{c}^{-} = 0.8/1 =
\qty{0.8}{\micro M}$; $C_{\text{ss}} = (1 + 0.8)/(5 + 1) = \qty{0.3}{\micro
M}$；$J = 5\times 0.3 - 1 = 0.5$ 亚基每秒，\qty{1.35}{nm/s}。
\end{solution}

\begin{solution}{exo:b3:cytoskeleton-motility:6}
$800/8 = 100$，即每秒 100 个 ATP。走完一米需要 $1/(8\times 10^{-7}) =
1.25\times 10^{6}$ 秒，约两周，共需 $1.25\times 10^{8}$ 步
与 ATP。若靠扩散则需 \num{16000} 年：
相差十亿倍。
\end{solution}

\begin{solution}{exo:b3:cytoskeleton-motility:7}
$F_{\max} = 8.3\times 10^{-20}/3.6\times 10^{-8} = \qty{2.3}{pN}$。同样的
能量分摊到更长步长上，力就更小——相当于更长的杠杆。
\omterm{def:b3:cytoskeleton-motility:motors}{驱动蛋白}（步长短，\qty{6}{pN}）适合重载；\omterm{def:b3:cytoskeleton-motility:motors}{肌球蛋白} V（步长
较长）单位 ATP 行进距离更远，适合快速、轻载的运输。
\end{solution}

\begin{solution}{exo:b3:cytoskeleton-motility:8}
(a) 紫杉醇使\omterm{def:b3:cytoskeleton-motility:filaments}{微管}凝固，纺锤体无法搜索、捕获并纠正错误，
有丝分裂在\omterm{def:b3:dna-repair:ddr}{检查点}停滞；爬行细胞仍能持续突出，
但失去长期极性。(b) 秋水仙碱去除\omterm{def:b3:cytoskeleton-motility:filaments}{微管}，
无纺锤体，有丝分裂停滞；爬行仍借助肌动蛋白继续，
但失去持续的方向。(c) 细胞松弛素封住带刺端：无
突出，无爬行；有丝分裂继续，但胞质分裂失败，
形成双核细胞。(d) 抑制 Rac：无\omterm{def:b3:cytoskeleton-motility:migration}{片状伪足}，因此无法爬行；
分裂基本不受影响。
\end{solution}

\begin{solution}{exo:b3:cytoskeleton-motility:9}
停留时间量级与半衰期相当，\qty{20}{s}（时间常数
$20/\ln 2 \approx \qty{29}{s}$）；不动分数为 \qty{10}{\%}。
边缘处的聚合须同时供给前进
与逆向流动：$1 + 1.5 = \qty{2.5}{\micro m/min}$。
\end{solution}

\begin{solution}{exo:b3:cytoskeleton-motility:10}
生长中的末端保有由 GTP-微管蛋白组成的帽；其后发生水解
产生受应力的 GDP-微管蛋白。帽的丢失——一种随机事件，其概率
随生长减慢而升高——释放应力，末端
迅速缩短，直至新帽形成。略高于\omterm{thm:b3:cytoskeleton-motility:treadmilling}{临界
浓度}时，每根\omterm{def:b3:cytoskeleton-motility:filaments}{微管}在任一时刻要么带帽并
生长，要么无帽并解聚，因此长度分化为一批长的
与一批即将消失的，而不是聚集在平均值附近。
细胞由此快速探索其体积，并获得
选择性稳定的能力——到达动粒或皮层位点的正端
被捕获并保留，其余则在
数分钟内回收：搜索与捕获。
\end{solution}

\begin{solution}{exo:b3:cytoskeleton-motility:11}
在此尺度下惯性可忽略，流体方程
具有时间可逆性：沿原路返回的运动（如桨的往返）
使物体回到出发处，无论每一划的速度如何。
旋转的螺旋永远不会沿原路返回，其运动
具有手性且连续，因此产生净推力。在 \qty{100}{Hz}
的转速下，每转一圈一千个质子，
每秒每个马达 $10^{5}$ 个质子。
\end{solution}

\begin{solution}{exo:b3:cytoskeleton-motility:12}
气道\omterm{def:b3:cytoskeleton-motility:cilia}{纤毛}无法摆动，黏液与细菌得不到清除，且
感染反复发生。精子鞭毛以相同\omterm{def:b3:cytoskeleton-motility:cilia}{轴丝}为基础，
不能运动，导致不育。胚胎的结节处运动\omterm{def:b3:cytoskeleton-motility:cilia}{纤毛}驱动
向左的流动，确定左–右轴；无流动时，侧别
随机决定，因此一半患者的器官发生左右颠倒。
\end{solution}

\begin{solution}{pb:b3:cytoskeleton-motility:1}
\textbf{1.} $\qty{300}{\micro m^3} = 3\times 10^{-13}$ 升；$\times
2\times 10^{-4}$ 摩尔/升 $= 6\times 10^{-17}$ 摩尔 $= 3.6\times 10^{7}$ 个
分子；$1.8\times 10^{7}$ 个位于细丝中。
\textbf{2.} $1.8\times 10^{7}\times \qty{2.7}{nm} = \qty{4.9}{cm}$ 长的细丝；
每根长 \qty{0.25}{\micro m}，约 $2\times 10^{5}$ 根
细丝。
\textbf{3.} 封端蛋白封住了大多数带刺端，profilin 与 thymosin-$\beta$4
共同结合单体，使真正游离的 ATP-肌动蛋白
接近\omterm{thm:b3:cytoskeleton-motility:treadmilling}{临界浓度}；生长只限于细胞
解除封端的那些末端。
\textbf{4.} $C_{\text{ss}} = 2.2/12.9 = \qty{0.17}{\micro M}$; $J = 11.6
\times 0.17 - 1.4 = 0.6$ 亚基每秒；\qty{1.6}{nm/s} $=
\qty{0.1}{\micro m/min}$.
\textbf{5.} $1/0.1 = 10$ 倍。
\textbf{6.} \qty{1}{\micro m/min} $= \qty{16.7}{nm/s} = 6.2$ 个亚基
每秒每根细丝；$\times 2\times 10^{5} = 1.2\times 10^{6}$ 个 ATP 每
秒，约占细胞周转量的 $0.1\,\%$。
\textbf{7.} $1/(8\times 10^{-7}) = 1.25\times 10^{6}$ 秒 $\approx 14.5$
天；$1.25\times 10^{8}$ 步与 ATP。
\textbf{8.} $L^{2}/2D$：$(1)^{2}/(2\times 10^{-12}) = 5\times 10^{11}$ 秒
$\approx \num{16000}$ 年（走完一米所需）；$(10^{-5})^{2}/(2\times
10^{-12}) = 50$ 秒对应 \qty{10}{\micro m}。扩散在
细胞体内已足够；更长距离则需要马达。
\textbf{9.} $5\times 10^{4}/6.02\times 10^{23} = 8.3\times 10^{-20}$ 焦耳
$= 20\,k_{B}T$.
\textbf{10.} $8.3\times 10^{-20}/8\times 10^{-9} = \qty{10}{pN}$;
效率 $6/10 \approx 60\,\%$。
\textbf{11.} $F = 6\pi\times 0.01\times 10^{-7}\times 8\times 10^{-7} =
1.5\times 10^{-14}$ 牛 $= \qty{0.015}{pN}$，比\omterm{prop:b3:cytoskeleton-motility:physics}{失速力}低四百倍：
单个马达足以克服黏性阻力；真正的
负载是障碍物与系绳。
\textbf{12.} $10^{4}\times 100 = 10^{6}$ 个 ATP 每秒，占神经元
能量预算的 $0.1\,\%$：运输代价低。但
马达在 ATP 供应下降时即停止，因此线粒体故障
会使突触缺乏细胞体中合成的一切物质——正是
神经退行性疾病中所见的轴突运输障碍。
\textbf{13.} \qty{10}{\micro m/min} $= \qty{167}{nm/s}$; $167/2.7 = 62$
亚基每秒每根细丝。
\textbf{14.} $200\times 20 = 4000$ 根细丝；$4000\times 62 = 2.5\times
10^{5}$ 亚基每秒。
\textbf{15.} $2.5\times 10^{5}$ 个 ATP 每秒，占
能量预算的 $0.025\,\%$。
\textbf{16.} $\qty{3}{\micro m}/(\qty{10}{\micro m/min}) = \qty{0.3}{min}
= \qty{18}{s}$；$2.5\times 10^{5}\times 18 = 4.5\times 10^{6}$ 个亚基
位于网络中。
\textbf{17.} $1\,\text{pN}\times 2.7\,\text{nm} = \qty{2.7}{pN.nm} =
0.66\,k_{B}T$：该能量的涨落出现的概率
约为 $e^{-0.66} \approx 0.5$——单亚基空隙持续开启，
棘轮机制完全合理。
\textbf{18.} 聚合本身即可推动，细菌充当其
自身网络前端的“膜”；借助细胞的
Arp2/3 机制，可达到\omterm{def:b3:cytoskeleton-motility:migration}{片状伪足}的速度，甚至更快——高达
\qty{0.5}{\micro m/s}, \qty{30}{\micro m/min}.
\textbf{19.} 当 $C = K_{d}$ 时，$\theta = 1/2$，且 $\mathrm{d}\theta/
\mathrm{d}C = 1/(4C)$：$C$ 相差 \qty{2}{\%} 时，$\Delta\theta
= 0.005$；每侧 \num{25000} 个受体，
前端多占据 $125$ 个。
\textbf{20.} 每侧约有 \num{12500} 个受体被占据，其涨落为
$\sqrt{\num{12500}} \approx 110$，因此两侧之差的涨落约为
$160$：$125$ 在任一
瞬间都淹没于噪声之中。受体约每秒重新结合一次；
对 $N$ 秒取平均，可使噪声减小 $\sqrt{N}$ 倍——十秒可将其
降至 $50$，梯度便可被读取。
\textbf{21.} 前端为 Rac（与 Cdc42 一起），后端为 Rho。
若 Rac 处处保持活性，细胞四周皆突出、铺展，却哪里也去不了。
\textbf{22.} $30/3 = \qty{10}{\micro m/min}$：与题中所给一致。
\textbf{23.} 突出被压力驱动的泡状突起所取代；黏附
与收缩仍依赖肌动蛋白和\omterm{def:b3:cytoskeleton-motility:motors}{肌球蛋白}，因此肌动蛋白在
循环的其余阶段仍不可或缺。
\textbf{24.} 前端每 \qty{18}{s} 重建一次：Rac 活性的改变
可在一个网络寿命内重定向聚合，新的
前端引导，后部的慢步骤随后跟进。
\textbf{25.} 体外\omterm{thm:b3:cytoskeleton-motility:treadmilling}{踏车行为}速率 \qty{1.6}{nm/s}；约 $14.5$ 天
穿过轴突所需；约 $2.5\times 10^{5}$ 个 ATP 每秒用于
突出。
\end{solution}
